\cgeentry{cge_closure}{Choose the model closure}

\cgesec{Description}

\cmd{cge_closure} sets the closure of the calibrated model: which variables
are fixed and which the model determines, for the external account, the
government, the savings--investment balance, the numeraire and each factor,
and, element by element, through swaps. It checks that the closure has as many
free variables as equations and that no equation is left without a free
variable, and displays the closure. Without options, it displays the current
closure.

\cgesec{Quick start}

Display the current closure
\begin{stcode}
cge_closure
\end{stcode}
A fixed exchange rate, investment-driven savings and capital fixed in each
activity
\begin{stcode}
cge_closure, external(fixexr) savinv(investment) factors(cap:specific)
\end{stcode}
Fixed government savings, with sales tax rates adjusting
\begin{stcode}
cge_closure, government(fixgsav sales)
\end{stcode}
Fix government savings and let real government consumption adjust
\begin{stcode}
cge_closure, swap(GSAV GADJ)
\end{stcode}
Remove the swaps
\begin{stcode}
cge_closure, swap(none)
\end{stcode}

\cgesec{Syntax}

\begin{stsyntax}
\cmdb{cge_closure} [\texttt{,} \argu{options}]
\end{stsyntax}

\begin{opttable}
\optgroup{Main}
\opt{ext}{ernal}\cmd{(flexexr|fixexr)} & external account; default \cmd{flexexr}\\
\optarg{gov}{ernment}{spec} & government: \cmd{flexgsav}, the default, or \cmd{fixgsav} with the adjusting tax\\
\opt{sav}{inv}\cmd{(savings|investment)} & savings--investment balance; default \cmd{savings}\\
\opt{num}{eraire}\cmd{(cpi|dpi)} & numeraire; default \cmd{cpi}\\
\optarg{fac}{tors}{spec} & closure of each factor, as \argu{factor}\cmd{:}\argu{type}; default \cmd{mobile}\\
\optgroup{Swaps}
\optarg{}{swap}{fix free \textup{[}fix free \dots\textup{]}} & fix and free variable elements in pairs; \cmd{swap(none)} removes the swaps\\
\end{opttable}

The defaults are those that \manref{cge_calibrate} sets. An option that is not
specified keeps its current setting.

\cgesec{Options}

\begin{optdesc}
\item[\opt{ext}{ernal}\cmd{(flexexr|fixexr)}] sets the closure of the
external account. \cmd{flexexr} fixes foreign savings, $\V{FSAV}$, in foreign
currency, and the exchange rate, $\V{EXR}$, adjusts; \cmd{fixexr} fixes the
exchange rate, and foreign savings adjust.

\item[\optarg{gov}{ernment}{spec}] sets the closure of the government.
\cmd{flexgsav} fixes tax rates and real government consumption, and
government savings, $\V{GSAV}$, adjust. \cmd{fixgsav} fixes government savings
and scales the rates of one tax in proportion to keep them: direct tax rates
(\cmd{fixgsav} or \cmd{fixgsav direct}, scaled by $\V{TINSADJ}$), sales tax
rates (\cmd{fixgsav sales}, $\V{TQADJ}$), import tariffs
(\cmd{fixgsav import}, $\V{TMADJ}$) or VAT rates (\cmd{fixgsav vat},
$\V{TVADJ}$). The adjusting tax must raise revenue in the base; rates that are
zero in the base stay zero. A colon may replace the space, as in
\cmd{fixgsav:vat}, which is how the closure is displayed. With real government
consumption fixed, \cmd{fixgsav} keeps a policy change revenue neutral: the
scaling factor shows the rise in the rates of the adjusting tax that replaces
the revenue lost.

\item[\opt{sav}{inv}\cmd{(savings|investment)}] sets the
savings--investment closure. \cmd{savings} fixes the savings rates, scaled by
$\V{MPSADJ}$, and investment adjusts to the savings available, through
$\V{IADJ}$; \cmd{investment} fixes real investment, and the savings rates of
the domestic institutions are scaled to finance it.

\item[\opt{num}{eraire}\cmd{(cpi|dpi)}] chooses the numeraire, the
price index fixed in the closure: \cmd{cpi}, the consumer price index, or
\cmd{dpi}, the producer price index of domestic sales.

\item[\optarg{fac}{tors}{spec}] sets the closure of each factor, with
\argu{spec} a list of \argu{factor}\cmd{:}\argu{type}. Factors that are not
listed are \cmd{mobile}, whatever their earlier closure. The types are
\cmd{mobile}, fully employed and mobile across activities, with a fixed supply
and one average price; \cmd{specific}, fixed in each activity, which pays its
own price; \cmd{unemployed}, a fixed price, with the supply adjusting;
\cmd{wagecurve}, under which the real wage falls as the unemployment rate
rises; and \cmd{minwage}, under which the real wage cannot fall below a floor:
either the floor binds and unemployment absorbs the shock, or the wage is
above the floor and the labour force is fully employed, and the solver finds
which. \cmd{wagecurve} and \cmd{minwage} need a labour force, from
\manref{cge_calibrate} with \opt{}{unemployment()} or \opt{}{minwage()};
\cmd{wagecurve} also needs the elasticity of \opt{}{wagecurve()}. Under both,
the wage and employment are free and the labour force is fixed
(chapter~\ref{chap:labour}).

\item[\optarg{}{swap}{fix free \textup{[}fix free \dots\textup{]}}] changes the
closure element by element, after the options above. Each pair fixes the first
element, which the closure leaves free, and frees the second, which it fixes.
Elements are written \argu{NAME}, \argu{NAME}\cmd{(}\argu{e1}\cmd{)} or
\argu{NAME}\cmd{(}\argu{e1}\cmd{,}\argu{e2}\cmd{)}, for example \cmd{W(lab)};
each element may appear once. Swaps replace any earlier ones and stay until
changed; \cmd{swap(none)} removes them. Variables that the reduced system
computes from their definitions, such as household demand, cannot be fixed.
\end{optdesc}

\cgesec{Remarks and examples}

The model has more variables than equations. The closure chooses which
variables are fixed, at their current values, and which the model determines.
It states how the economy is assumed to adjust to a shock, and results can
depend on it as much as on the shock. \manref{cge_calibrate} sets the default
closure, with no swaps; the closure set by \cmd{cge_closure} then stays until it
is changed or the model is calibrated again. \manref{cge_solve} solves under
the current closure and stores it with each simulation; it is saved with the
model by \manref{cge_save}, kept in sensitivity analysis and written into the
GAMS program by \manref{cge_gams}. \Chapref{chap:closures} lists the variables
that each option fixes and frees.

In the examples, the synthetic SAM is mapped and calibrated as in
\manref{cge_solve}.

\cgesubsec{Example 1: the default closure}

Without options, \cmd{cge_closure} displays the closure. After calibration it
is the default:
\stlog{clos_default}
\res{r(closure)} gives the closure in one line, as it is stored with each
simulation.

\cgesubsec{Example 2: other closure options}

A fixed exchange rate, investment fixed in real terms and capital fixed in
each activity:
\stlog{clos_alt}
The government and numeraire closures, which were not specified, keep their
settings, as does labour, which \opt{fac}{tors()} does not list and which is
therefore mobile.

\cgesubsec{Example 3: revenue-neutral tax replacement}

The other blocks are first returned to their defaults. With government savings
fixed and sales tax rates adjusting, the revenue lost when import tariffs are
removed is replaced by higher sales taxes:
\stlog{clos_fixgsav}
Government savings stay at their base value, direct tax rates and real
government consumption are unchanged, and the scaling factor of the sales tax
rates, \cmd{TQADJ}, rises by about 29 per cent:
\stlog{clos_fixgsav_res}
Under such a closure a shock to the rates of the adjusting tax itself, such as
removing tariffs while tariffs adjust, is not meaningful: the scaling factor
offsets it.

\cgesubsec{Example 4: swaps}

Government purchases of services can be fixed, and so shocked, by freeing the
scaling factor of government consumption, $\V{GADJ}$:
\stlog{clos_swap}
The swap is displayed with the closure. Government purchases are
$\V{QG}_c = \overline{\V{qg}}_c\,\V{GADJ}$ \eqref{eq:I7}, so the other purchases
rise in the same proportion:
\stlog{clos_swap_res}
Other common swaps are \cmd{swap(GSAV GADJ)}, which fixes government savings
and lets real government consumption adjust, and \cmd{swap(EXR CPI)}, which
makes the exchange rate the numeraire. A user variable defined by
\manref{cge_variable} can also be swapped into the closure, which fixes the
expression that defines it.

\cgesubsec{Example 5: closures that are rejected}

A closure is checked when it is set. Fixing government purchases of services
while freeing the direct tax scaling factor leaves the number of free variables
unchanged, but the equation of government demand for services then has no free
variable:
\stlog{clos_singular}
Such a closure makes the Jacobian singular, and \cmd{cge_closure} names the
equation. A swap must also fix a variable that is free and free one that is
fixed:
\stlog{clos_fixed}
After an error the closure in memory is unchanged; here the swap of
Example~4 was still in force until \cmd{swap(none)} removed it.

\cgesubsec{Closures, constraints and user equations}

Each constraint set by \manref{cge_constraint} adds an equation and solves for
a variable that the closure fixes, so the system stays square; the limited
variable must be free under the closure and the adjusting variable fixed, and
every later closure is checked against the constraints. The rents
$\V{RM}_c$, $\V{RE}_c$ and $\V{RA}_a$ are fixed at zero by every closure unless
a constraint frees them. A user variable given a value by \manref{cge_variable}
is free and needs a user equation of its own (\manref{cge_equation}).

\cgesec{Stored results}

\cmd{cge_closure} stores the following in \cmd{r()}:
\begin{storedres}
\resgroup{Scalars}
\res{r(N_swaps)} & number of swaps\\
\resgroup{Macros}
\res{r(closure)} & closure in one line, with the swaps\\
\end{storedres}

\cgesec{Methods and formulas}

Let $x$ be the vector of every element of every variable that exists in the
model, $\mathcal{F}$ the set of elements that the closure fixes, $N$ the number
of equation elements and $K$ the number of constraints. The closure options
build $\mathcal{F}$ block by block; each swap $(y_k, x_k)$ then requires
$y_k \notin \mathcal{F}$ and $x_k \in \mathcal{F}$, and replaces $\mathcal{F}$ by
$(\mathcal{F} \cup \{y_k\}) \setminus \{x_k\}$. The closure is accepted only if
it is square,
\begin{equation*}
\#\{\, j : x_j \text{ exists},\ j \notin \mathcal{F} \,\} + K = N,
\end{equation*}
where the adjusting variable of each constraint, although fixed, is solved
for. It must also leave every equation with at least one free variable, over
the pattern of the Jacobian (the variables that each equation element
contains); the equations with exact derivatives, which contain many
variables, are not part of this check. Swaps leave the count of free
variables unchanged. A square closure that passes the check can still be
singular at a point; \manref{cge_solve} then reports the smallest pivot of
the Jacobian.

\Chapref{chap:closures} gives the variables that each option fixes and
frees (\secref{sec:closure-options}) and the rules for swaps
(\secref{sec:swaps}).

\cgesec{Also see}

\begin{alsosee}
\seeentry{cge_shock}{Change a parameter or a fixed variable}
\seeentry{cge_solve}{Solve the model}
\seeentry{cge_constraint}{Limits that bind in some scenarios and not in others}
\seeentry{cge_variable}{User variables, defined by an expression or determined by user equations}
\seeentry{cge_calibrate}{Calibrate the model to the SAM}
\seechapter{chap:closures}{Closures}
\seechapter{chap:labour}{Labour-market regimes}
\end{alsosee}
